\section{Research Agenda and Conclusion}

\subsection{Research Agenda}

Mise en place as described here rests on a single case study and conceptual
argument. Empirical validation across practitioners, domains, and agent
models is essential. We propose five research questions to guide this work:

\begin{description}
  \item[RQ1:] \emph{Does deliberate preparation reduce agent misalignment
    compared to iterative development, and at what cost-benefit threshold?}
    Controlled experiments comparing preparation-first and iteration-first
    workflows on equivalent tasks could establish whether the upfront time
    investment yields net savings---and identify the project characteristics
    where each approach dominates.

  \item[RQ2:] \emph{What constitutes ``sufficient'' context for agentic
    work---can we define stopping criteria for the preparation phase?}
    A saturation-style analysis of context documents, measuring alignment
    quality as a function of context volume and specificity, would help
    practitioners know when to stop preparing and start building.

  \item[RQ3:] \emph{How does context fluency vary across practitioners,
    domains, and agent models?} Studies comparing developers with different
    backgrounds (e.g., education, domain expertise, years of experience) on
    identical agentic tasks could reveal whether context fluency is a
    trainable skill or a function of prior domain
    knowledge~\cite{polanyi1958personal}.

  \item[RQ4:] \emph{Can mise en place principles scale beyond prototyping to
    sustained, multi-month development projects?} Our hackathon case
    represents a constrained, time-limited setting. Longitudinal case studies
    tracking preparation practices and their outcomes across full product
    development cycles would test the methodology's applicability at
    scale~\cite{dakhel2025humanai}.

  \item[RQ5:] \emph{What is the relationship between preparation time and
    implementation quality across different project complexities?} A
    multi-project dataset correlating preparation artifacts (document count,
    word count, task granularity) with implementation metrics (defect rate,
    alignment accuracy, rework ratio) could yield quantitative models of the
    preparation-quality relationship~\cite{alomar2025promptware}.
\end{description}

These questions span experimental, observational, and longitudinal designs.
We believe the most immediately tractable is RQ1, which could be addressed
with paired tasks and a modest participant pool.

\subsection{Conclusion}

Mise en place is not a finished methodology but a starting framework---a
formalization of preparation practices that experienced practitioners are
already converging on independently~\cite{huntley2025ralph, horthy2025rpi,
yegge2025beads}. Our single case study suggests that deliberate preparation
can dramatically compress implementation time and reduce architectural
misalignment, but we make no claim of generalizability from one hackathon.
What we do claim is that the alignment problem in agentic coding is
fundamentally a preparation problem, and that the field would benefit from
treating the work that happens \emph{before} agents write code as a
first-class engineering discipline.

The concept of context fluency offers a lens for understanding emerging
practitioner skills in AI-assisted development---skills that bridge domain
expertise, pedagogical design, and specification craft. We invite empirical
validation, replication across domains and practitioner profiles, and the
development of tools that support preparation-phase
work~\cite{anthropic2025context, github2025speckit}. The code, increasingly,
is the easy part. The hard part---knowing what to build and
why---remains a distinctly human contribution.
